% !TeX root = ../../Book3.tex
\section{\texorpdfstring{\(\zeta\)}{\AnalysisPdfZeta} 函数与素数}
\label{b3:sec:2702}

带非负系数的对数导数把 \(\psi\) 转成一个解析函数。极点将提供主项，边界零点则可能产生与主项同量级的振荡，所以二者要分别处理。

\subsection{对数导数}
\label{b3:subsec:270201}

前章与本章采用同一个等式
\[
F(s)\coloneqq-\zeta'(s)/\zeta(s)
=\sum_{n\ge1}\Lambda(n)n^{-s}
=s\int_1^\infty\psi(x)x^{-s-1}\,\mathrm{d}x,
\quad\operatorname{Re}s>1.
\]
最后一步是 Abel 求和，端点 \(\psi(x)x^{-s}\to0\) 由 \(\psi(x)=O(x)\) 保证。由于 \(\zeta(s)=1/(s-1)+\gamma+O(s-1)\)，有 \(F(s)=1/(s-1)-\gamma+O(s-1)\)。

\subsection{Euler 乘积}
\label{b3:subsec:270202}

在 \(\sigma>1\)，Euler 乘积给
\[
\operatorname{Re}F(\sigma+\mathrm{i}t)
=\sum_{p,r\ge1}(\log p)p^{-r\sigma}\cos(rt\log p).
\]
单个余弦项会改变符号；把三个高度的值按合适系数组合后，则可使每项非负。绝对收敛允许先在每个素数幂处组合，再求全和。

\subsection{显式公式观点}
\label{b3:subsec:270203}

从形式上看，\(F(s)x^s/s\) 在 \(s=1\) 的留数为 \(x\)，而在零点 \(\rho\) 处的留数为 \(-m_\rho x^\rho/\rho\)。这说明零点实部为什么影响误差幂次。本章的证明采用 Tauber 方法，无需先建立一般反演显式公式；这段留数计算只解释单个奇点的局部贡献，不将无穷多留数未经截断相加。
